\section{The Target Compound and Related Structures}
\label{sec:identity}

\subsection{Scope of the Question}

Before synthesis conditions can be compared, the object of study has to be pinned down. The available sources are the original report of the Zn-containing target phase, historical studies of the Zn-free tetragonal Ba--Y silicates, and structurally related compounds that offer nothing more than a structural reference. The original report writes the target phase as $\mathrm{Ba_5Y_{12}Zn[O(SiO_4)]_8}$ and determines its structure by single-crystal X-ray diffraction~\cite{ababaikeri2024ba5y12zn}; the other two groups come close to the target phase in composition or crystallographic description but cannot stand in for the Zn-containing phase itself. This section checks, in turn, the original formula and its source and the revision of the historical nomenclature by modern structural work, and then states which conclusions the existing structural and phase-region evidence can support; the synthesis variables themselves are given in Section~\ref{sec:condition-matrix}.

This keeps the question ``which phase was obtained'' separate from the question ``which conditions were used''. For each record the table retains the formula as written in the original, the typeset form used here, the structural characterization result, and the open questions. Only a study in which both the formula and direct structural evidence point to the target phase counts as direct structural evidence for it. Sources on the same tetragonal Ba--Y silicate family are used to spot misassignment, non-stoichiometry, and competing phases; structurally related compounds merely flag local structural units that need checking and do not change the definition of the target phase.

\subsection{Provenance of the Target Phase}

Among the publications checked here, $\mathrm{Ba_5Y_{12}Zn[O(SiO_4)]_8}$ has a single direct report~\cite{ababaikeri2024ba5y12zn}. In that work the crystals were grown from a high-temperature solution in an open system and the structure was identified by single-crystal X-ray diffraction~\cite{ababaikeri2024ba5y12zn}. It remains the only study that confirms the target phase directly, and its results cannot be merged with the conditions of related compounds into a ``consensus recipe''. The report supports the target formula, the route type, and the single-crystal structure assignment, but it does not demonstrate cross-laboratory reproduction or a bulk phase-pure window, and it leaves the oxygen-stoichiometry expression, the Zn site, and the relationship to the Zn-free tetragonal silicate family only partly resolved.

Table~\ref{tab:identity-ledger} lists the original spelling and the typeset form used here in separate columns. The typeset form only harmonizes symbols, subscripts, and brackets; it neither rewrites variable-composition formulas as fixed-composition ones nor alters a published formula on the grounds of element balance.

\begin{table}[bp]
  \centering
  \caption{\textbf{Compositional and structural evidence for the target compound and related Ba--Y silicates.} Records are grouped by how the subject of each publication relates to the target phase; neither the Zn-free studies nor the structurally related studies constitute an independent reproduction of the Zn-containing target phase.}
  \label{tab:identity-ledger}
  \begin{tabular}{P{\dimexpr 0.1865\textwidth-2\tabcolsep\relax}P{\dimexpr 0.0813\textwidth-2\tabcolsep\relax}P{\dimexpr 0.2043\textwidth-2\tabcolsep\relax}P{\dimexpr 0.1393\textwidth-2\tabcolsep\relax}P{\dimexpr 0.3885\textwidth-2\tabcolsep\relax}}
    \toprule
    \thead{Formula (original \textrightarrow{} as typeset here)} & \thead{Ref.} & \thead{Main structural evidence} & \thead{Relation to target phase} & \thead{Established and open}\\
    \midrule
    $\mathrm{Ba_5Y_{12}Zn}$\allowbreak $\mathrm{[O(SiO_4)]_8}$ & \cite{ababaikeri2024ba5y12zn} & Single-crystal X-ray structure determination~\cite{ababaikeri2024ba5y12zn} & Direct report of the target compound & Established: Zn-containing target phase and its route type; open: independent reproduction, bulk phase-pure window, correspondence to the Zn-free family\\
    \addlinespace[3pt]
    $\mathrm{Ba_{5+x}Y_{13}Si_8O_{41}}$ (tentative; $x$ kept variable here) & \cite{kolitsch2009crystal,wierzbickawieczorek2017high} & Reported as an accompanying phase; systematic phase-formation screening~\cite{kolitsch2009crystal,wierzbickawieczorek2017high} & Same-family tetragonal Ba--Y silicate & Established: historical name and composition-screening results; open: exact stoichiometry, structural model, relation to the target phase\\
    \addlinespace[3pt]
    $\mathrm{Ba_xY_{26}Si_{16}}$\allowbreak $\mathrm{O_{71+x}}$ (non-stoichiometric formula kept here) & \cite{yamane2024synthesis} & Single-crystal revision to a non-stoichiometric modulated structure~\cite{yamane2024synthesis} & Same-family tetragonal Ba--Y silicate & Established: current structural reading of the historical formula; open: whether Zn can enter this family\\
    \addlinespace[3pt]
    $\mathrm{Ba_xY_{26}Si_{16}}$\allowbreak $\mathrm{O_{71+x}}$ ceramics & \cite{yamane2024microstructure} & Bulk phase composition and microstructure~\cite{yamane2024microstructure} & Same-family tetragonal Ba--Y silicate & Established: bulk characterization requirements for the non-stoichiometric family; open: separating secondary phases and extraneous SiO$_2$ from the intrinsic phase region\\
    \addlinespace[3pt]
    $\mathrm{Ba_5Y_{13}}$\allowbreak $\mathrm{[SiO_4]_8O_{8.5}}$ & \cite{gulay2024navigation} & EDS, cRED, synchrotron PXRD, and neutron refinement~\cite{gulay2024navigation} & Same-family tetragonal Ba--Y silicate & Established: independently confirmed tetragonal Ba--Y structure and phase-region reference; open: site correspondence to the Zn-containing target phase\\
    \addlinespace[3pt]
    $\mathrm{BaY_{16}Si_4O_{33}}$ & \cite{motozawa2022bay16si4o33} & Single-crystal structure with isolated orthosilicate units~\cite{motozawa2022bay16si4o33} & Structurally related compound & Established: structural reference for Ba--Y orthosilicates; open: no evidence that it is the same phase as the Zn-containing target\\
    \bottomrule
  \end{tabular}
\end{table}

\subsection{Reassessment of the Tetragonal Silicate Family}

That a historical name is still in use does not mean the structure has carried over. An early flux-growth study recorded $\mathrm{Ba_{5+x}Y_{13}Si_8O_{41}}$ as a tentative accompanying phase~\cite{kolitsch2009crystal}, and a later systematic flux study folded it into a multivariable phase-formation screen, but the subject remained the Zn-free tetragonal Ba--Y silicate~\cite{wierzbickawieczorek2017high}. Modern single-crystal work then recast the related nominal composition as the non-stoichiometric modulated structure $\mathrm{Ba_xY_{26}Si_{16}O_{71+x}}$~\cite{yamane2024synthesis}, and the companion ceramic study showed that bulk samples must also be checked for secondary phases and for SiO$_2$ picked up from the grinding medium~\cite{yamane2024microstructure}. Historical names therefore have to be re-verified against structural models and bulk phase composition; none of these studies directly observed the exact Zn-containing target formula.

An independent phase-region study combining EDS, cRED, synchrotron PXRD, and neutron refinement confirmed tetragonal $\mathrm{Ba_5Y_{13}[SiO_4]_8O_{8.5}}$~\cite{gulay2024navigation}, giving the Zn-free Ba--Y silicates a structural reference that is independent of the tentative nomenclature. The single-crystal structure of $\mathrm{BaY_{16}Si_4O_{33}}$ provides a Ba--Y reference structure built from isolated orthosilicate units~\cite{motozawa2022bay16si4o33}; it can be used to check how the polyhedra connect to the silicate groups, but it is not direct evidence for the Zn-containing target phase. On this basis we fix the formula of the target phase; the remaining records describe related structures and competing phases, and the scope of the claims they can support is set out in Section~\ref{sec:evidence-distance}.
